% !TeX program = XeLaTeX
% !TeX root = ../pujavidhanam.tex
\chapt{सूर्य-अर्घ्य-प्रदानम्}

\centerline{\bfseries स्त्रीणां सूर्यार्घ्यप्रदान-विधिः}
\centerline{\sffamily \emph{Method of women offering Arghya to Sūrya}}

{\sffamily After bathing, wearing clean clothes, applying tilakam etc., and taking clean water in a tīrtha pātram, this must be offered with devotion.}

\faClockO{} {\sffamily This must be offered \textbf{thrice every day}. It must be done facing \textbf{east} in the morning, and facing \textbf{west} in the evening. At noon, face east until midday has elapsed, and west thereafter.}

{\sffamily Pray to Sūrya with the first shloka. Then, chanting the second shloka three times, offer Arghyam for each repetition. \textbf{Arghya} means joining the two hands together, cupped with palms facing up, taking water in them, and offering it with devotion.}

\centerline{\bfseries करणीयानि, अकरणीयानि च}
\centerline{\sffamily \emph{Do-s and don't-s}}

{\sffamily Arghya must be offered in a waterbody such as a river, or in a clean place. If that is not possible, it may be offered in a plate or vessel and then poured away in another clean place.

The water offered as Arghya should not get mixed with dirty water; it should not be touched by feet. If it spills where feet might touch it, wipe it away.

Not only for this, but for any anushthanam, the vessels used must be made of traditional metals such as copper, brass etc. Stainless steel should not be used.

If one incurs an āśaucham, one may mentally pray to Sūrya, but perform this only after the āśaucham has elapsed and śuddhi is acquired.}

\dnsub{प्रार्थना}

\twolineshloka*
{एहि सूर्य सहस्रांशो तेजोराशे जगत्पते}
{अनुकम्पय मां भक्त्या गृहाणार्घ्यं दिवाकर}

{\sffamily Come O Sūrya! You are the thousand-rayed effulgence! You are the ruler of this world! O causer of daytime! Please accept my Arghyam, showing compassion towards my devotion.}

\dnsub{अर्घ्य-प्रदानम्}

\twolineshloka*
{नमो विवस्वते ब्रह्मन् भास्वते विष्णुतेजसे}
{जगत्सवित्रे शुचये सवित्रे फलदायिने}
\centerline{(एवं त्रिः)}

{\sffamily O Supreme Being who shines with many rays! Your brilliance pervades everywhere! You are the one who created this world. You are pure. You are the witness to all the happenings of this world. You are the one who gives the fruits of actions to beings! To You my prostrations!}

\closesection
